\section{Philosophical lineage: \DCI\ and the nomological machine}
\label{sec:lineage}

The frame is not new metaphysics; it is the operationalization of an existing one. Cartwright's \emph{The Dappled World} \citep{cartwright1999} argues that laws of nature are not universal truths holding everywhere and always, but descriptions of the behavior of \textbf{nomological machines} --- stable arrangements of components with fixed capacities that, shielded from interference and set running repeatedly, produce the regular behavior we then codify as law. On this view the world is \emph{dappled}: a patchwork of local regularities, each valid within its machine, not a single seamless deductive system.

Two of this paper's load-bearing commitments are exactly Cartwright's, in computational dress.

First, \textbf{the executable rule substrate $\Rule$ is a nomological machine}. It fixes the components and their capacities, shields the dynamics from outside interference, and can be run arbitrarily many times to generate ground truth. What we called ``closed-loop discovery against a verifier'' is, in Cartwright's terms, an agent operating a nomological machine and reading off its lawful outputs. The verifier is not a convenience; it is the machine that makes lawfulness available at all.

Second, \textbf{\DCI's domain-boundedness is Cartwright's machine-locality, not a limitation to be apologized for}. \DCI\ does not aspire to universal law and then fall short; it aspires to \emph{completeness within a machine}, which on Cartwright's account is where lawfulness genuinely lives. On this view --- a philosophical position, offered as a divergence from the universalist reading rather than a refutation of it --- universality is not so much withheld as reframed: the question ``why isn't this general?'' presupposes that lawful structure is global, whereas a dappled world makes domain-completeness the natural target. This inverts the usual reading of ``narrow'': a domain-complete agent is not a general agent with its scope clipped, but the \emph{natural} epistemic agent for a dappled world. Nothing here forecloses the pursuit of universal agency; the claim is only that the domain-complete rung is real, prior, and where most near-term value sits.

\textbf{Where a domain begins and ends.} Machine-locality invites a fair challenge: if a human must hand-craft the rule substrate $\Rule$ and the token library, is the agent autonomous or merely an efficient search inside a human-built cage? The honest answer distinguishes \emph{provisioning} a machine from \emph{operating} it: supplying $\Rule$ is design, not discovery, and domain-completeness is a claim about what the agent reaches \emph{once} $\Rule$ is fixed, not about inventing $\Rule$ from nothing (the same division of labour that lets us call a scientist creative without crediting them with authoring physics). That division yields three usable boundary criteria, offered as criteria rather than theorems. \emph{Individuation}: a domain is the largest region over which a single $\Prin$ stays invariant and $\Rule$ stays executable; the domain \emph{edge} is exactly where $\Prin$ ceases to be conserved, so a domain cannot expand without bound unless a universal invariant exists --- which is just the universalist claim this paper declines. \emph{Non-triviality}: a domain must admit more than one law derivable from $\Rule$ --- a generative grammar, not a single fact --- or it collapses into a narrow-AI task with nothing left to discover; the interesting domains sit between these poles. \emph{Composition}: two \DCI\ agents whose domains overlap compose exactly where their invariants are jointly satisfied. Chemical physics, for instance, is the intersection region where both the chemical and the physical $\Prin$ hold; a composite agent is domain-complete over that intersection and no further, and where the two invariants conflict neither agent's guarantee survives. Formalizing these criteria --- an explicit calculus of domain intersection and the conditions under which two $\Prin$ compose --- is future work (\S\ref{sec:testable}).

This also sharpens what $\Prin$ is. A vertically organizing principle is the invariant \emph{of a machine} (or of a family of machines that share it), not a Theory of Everything --- the physics cousin being the Laughlin--Pines ``protectorate,'' a higher organizing principle whose consequences hold within its regime irrespective of the underlying microphysics \citep{laughlin2000}. The lineage is therefore clean: Cartwright's nomological machines and dappled lawfulness $\rightarrow$ vertically organizing principles as machine-local invariants \citep{frasch2026a} $\rightarrow$ \DCI\ as the artificial agent that is \emph{complete relative to a machine}. Positioning \DCI\ this way reframes the ``why not just push for general AGI?'' reflex at the level of philosophy of science, not only engineering: to the extent that much of the lawful structure worth discovering is machine-local, the agent matched to it is domain-complete, not universal.
